\chapter{Methodology}

1. Artifact Excavation (Data Collection)

Objective: Gather a diverse and representative sample of software architectures.

	•	Source Code Repositories: Collect architectures from open-source projects on platforms like GitHub, GitLab, and Bitbucket.
	•	Documentation: Obtain architectural diagrams, design documents, and technical specifications from project documentation.
	•	Interviews and Surveys: Conduct interviews with software architects and developers to gain insights into architectural decisions and rationale.
	•	Legacy Systems Analysis: Include older software systems to understand historical architectural practices.

Considerations:

	•	Ethical Compliance: Ensure all collected data respects intellectual property rights and privacy laws.
	•	Diversity: Aim for a variety of software types (e.g., web applications, mobile apps, embedded systems) and domains (e.g., healthcare, finance, education).

2. Classification and Cataloging

Objective: Organize the collected architectures into meaningful categories.

	•	Structural Classification:
	•	Monolithic vs. Microservices: Distinguish between tightly coupled and loosely coupled systems.
	•	Layered vs. Event-Driven: Categorize based on architectural patterns.
	•	Functional Classification:
	•	Domain-Based: Group architectures by industry or application domain.
	•	Technology Stack: Classify based on programming languages, frameworks, and platforms used.
	•	Temporal Classification:
	•	Era-Based: Organize architectures by the period they were developed (e.g., mainframe era, client-server era, cloud era).

Tools:

	•	Databases: Use relational or graph databases to store and query architectural data.
	•	Taxonomy Software: Employ tools like Protégé for ontology development.

3. Structural Analysis

Objective: Analyze the internal components and their relationships within each architecture.

	•	Component Identification: Break down architectures into their fundamental components (e.g., modules, services, layers).
	•	Relationship Mapping: Diagram the interactions and dependencies between components.
	•	Pattern Recognition: Identify common architectural patterns (e.g., MVC, RESTful services) and anti-patterns.
	•	Metric Calculation:
	•	Complexity Metrics: Evaluate cyclomatic complexity, coupling, and cohesion.
	•	Scalability and Performance Indicators: Assess load times, response rates, and throughput.

Methods:

	•	Static Code Analysis: Use tools like SonarQube or Understand.
	•	Visualization: Create UML diagrams or use visualization tools like Gephi.

4. Contextual Analysis

Objective: Understand the architectures within their broader socio-technical contexts.

	•	Historical Context: Investigate the technological landscape and limitations during the time of development.
	•	Organizational Factors: Examine how company culture, team structure, and project management methodologies influenced architectural decisions.
	•	Economic and Market Forces: Consider market demands, competition, and economic constraints.

Data Sources:

	•	Project Histories: Review version histories, commit messages, and project timelines.
	•	Industry Reports: Consult market analyses and technology trend reports.

5. Comparative Analysis

Objective: Compare architectures to identify trends, commonalities, and divergences.

	•	Cross-Artifact Comparison: Analyze similarities and differences across architectures within the same category.
	•	Temporal Trends: Observe how architectures have evolved over different technological eras.
	•	Best Practices Identification: Determine architectural features associated with successful projects.

Approach:

	•	Statistical Analysis: Use statistical methods to find correlations.
	•	Case Studies: Develop in-depth studies of particular architectures.

6. Interpretation and Hypothesis Formation

Objective: Derive insights and formulate definitions based on the analyses.

	•	Definition Refinement: Use findings to define what constitutes a software architecture.
	•	Boundary Setting: Clarify what is not considered part of software architecture (e.g., low-level code optimizations).
	•	Theory Development: Propose theories on how and why certain architectural styles emerge.

Validation:

	•	Peer Review: Present findings to other researchers for feedback.
	•	Expert Consultation: Engage with seasoned software architects for validation.

7. Documentation and Preservation

Objective: Ensure that all findings are thoroughly documented and preserved for future research.

	•	Archiving: Store all collected data, analyses, and interpretations in accessible repositories.
	•	Standardized Reporting: Use consistent formats for documentation to facilitate sharing and collaboration.

Tools:

	•	Digital Libraries: Use platforms like Zenodo or Figshare.
	•	Version Control: Employ Git for tracking changes in research documents.

8. Dissemination

Objective: Share the research outcomes with the broader community.

	•	Publications: Write research papers, articles, and reports.
	•	Conferences and Workshops: Present findings at industry and academic events.
	•	Open-Source Contributions: If applicable, contribute tools or datasets back to the community.

Applying Archaeological Principles

To effectively borrow from archaeology, consider the following principles:

	•	Stratigraphy: Just as archaeologists study layers of soil, examine layers of software (e.g., presentation, business logic, data access) to understand the buildup and interaction over time.
	•	Typology: Develop a classification system for architectures, akin to how artifacts are grouped by material, form, or function.
	•	Seriation: Arrange architectures in a sequence based on certain attributes to study chronological development.
	•	Contextualization: Always interpret architectures within the context of their environment, recognizing that external factors significantly influence their form and function.

Defining Software Architecture

Based on this method, you can define software architecture by:

	•	Structural Elements: Identifying the key components and their configurations.
	•	Relationships and Interactions: Understanding how components interact and depend on each other.
	•	Principles and Guidelines: Recognizing the underlying principles that guide architectural decisions.
	•	Constraints: Acknowledging limitations imposed by technology, resources, or requirements.

What Software Architecture Is Not

Through this approach, you can delineate that software architecture is not:

	•	Just Code: It’s more than lines of code; it’s the high-level structure.
	•	Single Design Patterns: While design patterns inform architecture, they are not architectures themselves.
	•	Implementation Details: Architecture focuses on the big picture rather than low-level implementation specifics.

Benefits of This Approach

	•	Comprehensive Understanding: Gain a holistic view of software architectures across different contexts.
	•	Historical Insight: Learn how architectures have evolved and why certain practices emerged.
	•	Best Practices Identification: Isolate architectural strategies that lead to successful software systems.
	•	Cross-Disciplinary Innovation: Enrich information systems research with methodologies from archaeology.

Final Thoughts

By treating software architectures as artifacts, you open up new avenues for analysis and understanding. This archaeological approach allows you to systematically define software architecture, appreciate its complexities, and distinguish it from related concepts. It fosters a deeper appreciation of the socio-techni



Reasoning Behind the Proposed Method

Overview:

By adopting archaeological research methods, we can systematically analyze software architectures as artifacts that embody historical, cultural, and technical information. This approach allows us to:

	•	Deconstruct software architectures to understand their components and relationships.
	•	Contextualize architectures within their socio-technical environments.
	•	Compare and classify architectures to identify patterns and evolutions.
	•	Define what constitutes software architecture by distinguishing it from other software elements.

1. Artifact Excavation (Data Collection)

Reasoning:

	•	Comprehensive Sampling: Like archaeologists excavating artifacts from various sites, collecting a wide range of software architectures ensures a representative sample for analysis.
	•	Multiple Sources: Gathering data from code repositories, documentation, and expert interviews provides a holistic view, capturing both technical details and human insights.
	•	Inclusion of Legacy Systems: Studying older architectures helps trace the evolution of architectural practices and understand historical contexts.

2. Classification and Cataloging

Reasoning:

	•	Organization: Classifying architectures aids in organizing the data, making it manageable and facilitating pattern recognition.
	•	Multiple Classification Dimensions: By categorizing architectures structurally, functionally, and temporally, we capture different facets that may influence architectural decisions.
	•	Use of Tools: Employing databases and taxonomy software enhances accuracy and efficiency in cataloging.

3. Structural Analysis

Reasoning:

	•	Component Breakdown: Identifying fundamental components reveals the building blocks of architectures.
	•	Relationship Mapping: Understanding interactions between components highlights dependencies and communication pathways.
	•	Pattern Recognition: Identifying common patterns and anti-patterns helps in understanding best practices and pitfalls.
	•	Metric Calculation: Quantitative metrics provide objective measures to compare architectures.

4. Contextual Analysis

Reasoning:

	•	Holistic Understanding: Contextual factors significantly influence architectural decisions; thus, analyzing these factors provides insights beyond technical aspects.
	•	Historical Context: Recognizing technological limitations and advancements of different eras explains why certain architectural choices were made.
	•	Organizational and Economic Factors: These factors can constrain or drive architectural innovation, affecting the final design.

5. Comparative Analysis

Reasoning:

	•	Trend Identification: Comparing architectures helps identify common trends, shifts in practices, and emerging patterns.
	•	Learning from Differences: Understanding divergences can reveal why certain architectures succeed or fail in specific contexts.
	•	Best Practices: Isolating features associated with successful projects aids in defining effective architectural principles.

6. Interpretation and Hypothesis Formation

Reasoning:

	•	Synthesis of Findings: Interpreting the data allows for the formulation of a precise definition of software architecture.
	•	Distinguishing Boundaries: Clearly defining what is not part of software architecture prevents ambiguity and scope creep.
	•	Theory Development: Proposing theories based on empirical data strengthens the academic rigor of the research.

7. Documentation and Preservation

Reasoning:

	•	Knowledge Retention: Thorough documentation ensures that findings are preserved for future reference and research continuity.
	•	Accessibility: Storing data in accessible repositories promotes transparency and facilitates peer review.
	•	Standardization: Consistent documentation practices enhance collaboration and interoperability among researchers.

8. Dissemination

Reasoning:

	•	Community Contribution: Sharing findings enriches the collective understanding and spurs further research.
	•	Feedback Loop: Presenting at conferences and publishing allows for critique and refinement of the research.
	•	Practical Impact: Open-source contributions can have immediate benefits for practitioners in the field.

Ensuring Completeness

After reviewing the proposed method, here are additional elements to consider to ensure the research is comprehensive:

A. Incorporating Ethnographic Methods

Reasoning:

	•	Human Factors: Software architecture is not only a technical construct but also a product of human collaboration and decision-making.
	•	Cultural Influences: Team dynamics, organizational culture, and individual expertise can shape architectural outcomes.

Implementation:

	•	Participant Observation: Engage with development teams to observe the architectural decision-making process.
	•	Narrative Analysis: Analyze stories and anecdotes shared by practitioners to gain insights into tacit knowledge.

B. Technological Evolution Analysis

Reasoning:

	•	Tool Influence: The availability and popularity of certain tools and technologies can influence architectural choices.
	•	Innovation Adoption: Studying how new technologies are adopted into architectures can reveal patterns of innovation diffusion.

Implementation:

	•	Technology Lifecycle Mapping: Chart the emergence and decline of technologies within architectures over time.
	•	Adoption Rate Analysis: Examine how quickly new technologies are incorporated into existing architectures.

C. Socio-Economic Impact Assessment

Reasoning:

	•	Economic Drivers: Cost considerations, market competition, and resource availability can dictate architectural decisions.
	•	Societal Needs: Social demands and user expectations can shape the functionality and structure of software systems.

Implementation:

	•	Cost-Benefit Analysis: Assess how economic factors influence architecture, such as trade-offs between performance and cost.
	•	User Demand Studies: Analyze how user requirements impact architectural complexity and design.

D. Cross-Disciplinary Integration

Reasoning:

	•	Enriched Perspectives: Integrating insights from fields like anthropology, sociology, and cognitive science can deepen the understanding of architecture as a socio-technical artifact.
	•	Holistic Analysis: Cross-disciplinary approaches help in examining the interplay between technology and human factors.

Implementation:

	•	Interdisciplinary Workshops: Collaborate with experts from other fields to broaden the analytical framework.
	•	Literature Review: Incorporate theories and models from other disciplines that are relevant to architectural analysis.

E. Addressing Limitations and Biases

Reasoning:

	•	Research Validity: Acknowledging and mitigating biases enhances the credibility of the research.
	•	Generalizability: Understanding the limitations ensures that conclusions are appropriately framed.

Implementation:

	•	Bias Identification: Reflect on potential biases in data selection and interpretation.
	•	Limitations Section: Explicitly state the research limitations in documentation and publications.

Final Thoughts

Summary:

	•	The proposed method is grounded in adapting archaeological principles to systematically study software architectures.
	•	Each step is designed to uncover different layers of understanding, from technical structures to socio-cultural contexts.
	•	The inclusion of additional considerations like ethnographic methods and cross-disciplinary integration enriches the research.

Conclusion:

By thoroughly examining software architectures as artifacts, we can develop a nuanced definition that captures their essence and distinguishes them from other software elements. This comprehensive approach ensures that we consider all relevant factors, leading to robust and insightful findings.

What Software Architecture Is and Is Not

Based on the Method:

	•	Is:
	•	A high-level abstraction of a software system’s structure.
	•	Defined by its components, their interactions, and the guiding principles behind their arrangement.
	•	Influenced by technical, organizational, and socio-economic factors.
	•	Is Not:
	•	Merely the source code or low-level implementation details.
	•	A single design pattern or technology stack.
	•	Static; it evolves with changing contexts and requirements.

Ensuring Nothing Was Missed

	•	Comprehensiveness: The method covers data collection, analysis, interpretation, and dissemination.
	•	Multi-Faceted Analysis: Includes technical, historical, organizational, and socio-economic perspectives.
	•	Additional Considerations: Ethnographic methods, technological evolution, socio-economic impacts, cross-disciplinary integration, and bias mitigation have been incorporated.

Confidence:

With these additions and the thorough reasoning provided, the research method is comprehensive and well-suited to define what software architecture is and is not by viewing it through the lens of archaeology.